\documentclass[12pt,a4paper]{article}
\usepackage[utf8]{inputenc}
\usepackage{graphicx}
\usepackage{xcolor}
\usepackage{listings}
\usepackage{geometry}
\usepackage{fancyhdr}
\usepackage{hyperref}
\usepackage{titlesec}
\usepackage{etoolbox}

\geometry{margin=1in}

% Colors
\definecolor{primary}{RGB}{26, 26, 46}
\definecolor{secondary}{RGB}{22, 33, 62}
\definecolor{accent}{RGB}{78, 205, 196}
\definecolor{light}{RGB}{248, 249, 250}
\definecolor{warning}{RGB}{255, 107, 107}

% Header/Footer
\pagestyle{fancy}
\fancyhf{}
\rhead{\thepage}
\lhead{\textbf{Lambda Log Analyser - Technical Documentation}}
\cfoot{\thepage}

% Section styling
\titleformat{\section}{\Large\bfseries\color{primary}}{}{0em}{}
\titleformat{\subsection}{\large\bfseries\color{secondary}}{}{0em}{}
\titleformat{\subsubsection}{\normalsize\bfseries\color{secondary}}{}{0em}{}

\begin{document}

% Title Page
\begin{titlepage}
    \centering
    \vspace*{2cm}
    
    {\Huge\bfseries\color{primary} Lambda Log Analyser\par}
    \vspace{0.5cm}
    {\LARGE\color{secondary} Technical Documentation\par}
    \vspace{2cm}
    
    {\Large\bfseries Architecture, Features \& Process Pipeline\par}
    \vspace{1cm}
    
    \rule{0.8\textwidth}{0.5pt}
    
    \vspace{2cm}
    {\large Version 1.0\par}
    \vspace{1cm}
    {\large\ today\par}
    
    \vfill
    {\itshape\color{gray} A comprehensive technical guide to building, deploying, and using the Lambda Log Analyser dashboard.\par}
\end{titlepage}

\pagebreak
\tableofcontents
\pagebreak

% ==============================================================================
% SECTION 1: EXECUTIVE SUMMARY
% ==============================================================================
\section{Executive Summary}

The Lambda Log Analyser is a sophisticated dashboard application designed to analyze AWS Lambda function logs, extract meaningful performance metrics, detect errors, and provide actionable insights through an intuitive web interface.

\bigskip
\noindent\textbf{Key Capabilities:}
\begin{itemize}
    \item Real-time CloudWatch Logs fetching with automatic pagination
    \item Comprehensive error detection and classification
    \item Performance metrics extraction (duration, memory, cold starts)
    \item Cost estimation based on AWS Lambda pricing model
    \item Health score calculation for function monitoring
    \item Interactive visualizations with Plotly charts
    \item PDF and CSV export capabilities
\end{itemize}

% ==============================================================================
% SECTION 2: SYSTEM ARCHITECTURE
% ==============================================================================
\section{System Architecture}

\subsection{High-Level Architecture}

\begin{figure}[htbp]
    \centering
    \begin{verbatim}
    ┌─────────────────────────────────────────────────────────────────┐
    │                        USER INTERFACE                            │
    │                   (Streamlit Web Dashboard)                      │
    │                                                                 │
    │  ┌─────────────┐  ┌──────────────┐  ┌────────────────────────┐ │
    │  │ AWS Config  │  │ Log Groups   │  │ Analysis Controls      │ │
    │  └─────────────┘  └──────────────┘  └────────────────────────┘ │
    └─────────────────────────────────────────────────────────────────┘
                                   │
                                   ▼
    ┌─────────────────────────────────────────────────────────────────┐
    │                      APPLICATION LAYER                          │
    │                                                                 │
    │  ┌────────────────┐  ┌────────────────┐  ┌──────────────────┐  │
    │  │ LogFetcher     │  │ LogParser      │  │ AnalysisEngine   │  │
    │  │ (boto3)       │  │ (Regex-based)  │  │ (Metrics Calc)   │  │
    │  └────────────────┘  └────────────────┘  └──────────────────┘  │
    │                                                                 │
    │  ┌────────────────┐  ┌────────────────┐  ┌──────────────────┐  │
    │  │ ChartGenerator │  │ PDFExporter    │  │ CSVExporter      │  │
    │  │ (Plotly)      │  │ (ReportLab)    │  │ (Pandas)         │  │
    │  └────────────────┘  └────────────────┘  └──────────────────┘  │
    └─────────────────────────────────────────────────────────────────┘
                                   │
                                   ▼
    ┌─────────────────────────────────────────────────────────────────┐
    │                    AWS CLOUD SERVICES                           │
    │                                                                 │
    │  ┌────────────────────────────────────────────────────────┐    │
    │  │              Amazon CloudWatch Logs                    │    │
    │  │         (Log Groups, Log Streams, Events)               │    │
    │  └────────────────────────────────────────────────────────┘    │
    └─────────────────────────────────────────────────────────────────┘
    \end{verbatim}
    \caption{System Architecture Diagram}
\end{figure}

\subsection{Component Architecture}

\subsubsection{Streamlit Dashboard (app.py)}
The Streamlit application serves as the user interface layer, handling:
\begin{itemize}
    \item User authentication input (AWS credentials)
    \item Configuration management (region, log groups, lookback period)
    \item Session state management for analysis results
    \item Tabbed interface for different analysis views
    \item Real-time chart rendering with Plotly
    \item Export functionality (PDF, CSV, JSON)
\end{itemize}

\subsubsection{Log Fetcher (log\_fetcher.py)}
Responsible for CloudWatch Logs interaction:
\begin{itemize}
    \item AWS boto3 client initialization with configurable region
    \item Automatic pagination handling for large log streams
    \item Time-window filtering for lookback period
    \item Multi-log-group support
    \item Error handling for AWS API failures
\end{itemize}

\subsubsection{Log Parser (log\_parser.py)}
Core parsing engine using regex patterns:
\begin{itemize}
    \item REPORT line parsing for performance metrics
    \item START/END event detection
    \item Error pattern matching (multiple regex patterns)
    \item Cold start detection via Init Duration field
    \item Exception type extraction
\end{itemize}

\pagebreak

% ==============================================================================
% SECTION 3: DATA FLOW PIPELINE
% ==============================================================================
\section{Data Flow Pipeline}

\subsection{Pipeline Overview}

\begin{figure}[htbp]
    \centering
    \begin{verbatim}
    ┌──────────────┐    ┌──────────────┐    ┌──────────────┐    ┌──────────────┐
    │  User Input  │───▶│ Log Fetcher  │───▶│  Log Parser  │───▶│  Analysis    │
    │  (Config)    │    │  (boto3)     │    │  (Regex)     │    │  Summary     │
    └──────────────┘    └──────────────┘    └──────────────┘    └──────────────┘
                                                              │
                         ┌────────────────────────────────────┘
                         │
                         ▼
    ┌──────────────┐    ┌──────────────┐    ┌──────────────┐
    │  PDF Report  │◀───│  Visualize   │◀───│  Metrics     │
    │  Generator   │    │  (Charts)    │    │  Calculator  │
    └──────────────┘    └──────────────┘    └──────────────┘
    \end{verbatim}
    \caption{Data Processing Pipeline}
\end{figure}

\subsection{Step-by-Step Process}

\subsubsection{Step 1: Configuration}
User provides AWS credentials and configuration:
\begin{itemize}
    \item AWS Access Key ID and Secret Access Key
    \item AWS Region selection
    \item CloudWatch Log Group names (one per line)
    \item Lookback period in minutes
    \item Lambda memory size and timeout for cost estimation
\end{itemize}

\subsubsection{Step 2: Log Fetching}
The LogFetcher class performs:
\begin{enumerate}
    \item Initializes boto3 CloudWatch Logs client with provided credentials
    \item For each log group, lists active log streams
    \item Filters streams by last event timestamp (lookback window)
    \item Fetches events from each stream with pagination
    \item Yields raw event dictionaries with structure:
\end{enumerate}

\begin{lstlisting}[language=Python, caption=Raw Event Structure]
{
    "log_group": "/aws/lambda/function-name",
    "log_stream": "2024/01/01/[$LATEST]abc123",
    "timestamp": 1704067200000,  # Epoch milliseconds
    "message": "REPORT RequestId: abc-123 Duration: 312.45 ms ..."
}
\end{lstlisting}

\subsubsection{Step 3: Log Parsing}
The LogParser processes each raw event:

\bigskip
\noindent\textbf{REPORT Event Parsing:}
\begin{itemize}
    \item Regex pattern extracts: RequestId, Duration, Billed Duration, Memory Size, Max Memory Used, Init Duration
    \item Cold start detection: \texttt{is\_cold\_start = init\_duration is not None}
    \item Creates ParsedEvent with all extracted metrics
\end{itemize}

\bigskip
\noindent\textbf{Error Detection Patterns:}
\begin{itemize}
    \item \texttt{[ERROR]} or \texttt{ERROR:} in message
    \item Python Traceback blocks
    \item Exception types (ValueError, TypeError, etc.)
    \item Timeout messages
    \item Runtime errors (ExitError, OutOfMemoryError)
\end{itemize}

\subsubsection{Step 4: Analysis Summary}
The build\_summary() method aggregates:
\begin{itemize}
    \item Total invocations and errors
    \item Cold start count and rate
    \item Duration statistics (min, max, avg, P95, P99)
    \item Memory statistics (min, max, avg, efficiency)
    \item Billed duration for cost calculation
    \item Error events list with full messages
\end{itemize}

\pagebreak

% ==============================================================================
% SECTION 4: FEATURE SPECIFICATIONS
% ==============================================================================
\section{Feature Specifications}

\subsection{Overview Tab Features}

\subsubsection{Executive Summary Cards}
Five metric cards displaying:
\begin{itemize}
    \item \textbf{Total Events} - All processed log events with invocation delta
    \item \textbf{Errors} - Error count with percentage rate and color-coded delta
    \item \textbf{Cold Starts} - Count with rate percentage
    \item \textbf{Avg Duration} - Average execution time in milliseconds
    \item \textbf{Est. Cost} - Calculated AWS Lambda cost (6 decimal precision)
\end{itemize}

\subsubsection{Health Score Gauge}
Composite health indicator (0-100\%) calculated from:
\begin{itemize}
    \item Error rate (50\% weight)
    \item Cold start rate (10\% weight)
    \item P95 duration threshold (20\% weight)
    \item Memory efficiency (20\% weight)
\end{itemize}

Formula: \texttt{Health = 100 - (ErrorRate * 50) - (ColdStartRate * 10) - DurationPenalty - MemoryPenalty}

\subsubsection{Event Distribution Pie Chart}
Interactive Plotly pie chart showing:
\begin{itemize}
    \item Report events (green - successful invocations)
    \item Error events (red - failed invocations)
    \item Start/End events (other colors)
\end{itemize}

\subsection{Performance Tab Features}

\subsubsection{Duration Distribution Histogram}
Plotly histogram with:
\begin{itemize}
    \item 30 bins for duration distribution
    \item Average duration vertical line
    \item Interactive hover tooltips
    \item Export capability
\end{itemize}

\subsubsection{Memory Usage Histogram}
Similar to duration with memory values in MB.

\subsubsection{Invocation Timeline Chart}
Time-series scatter plot showing:
\begin{itemize}
    \item Duration over time (line chart)
    \item Cold start markers (star markers in red)
    \item Interactive zoom and pan
    \item Unified hover mode
\end{itemize}

\subsection{Cost Estimation}

\subsubsection{AWS Lambda Pricing Model}
\begin{verbatim}
Compute Cost = TotalBilledSeconds × MemoryGB × $0.0000166667/GB-second
Request Cost = TotalInvocations × $0.0000002/request
Total Cost = Compute Cost + Request Cost
\end{verbatim}

Example: 256 MB, 100 invocations, 30000 ms total billed:
\begin{itemize}
    \item Memory GB = 0.25
    \item Compute seconds = 30
    \item Compute cost = 30 × 0.25 × \$0.0000166667 = \$0.000125
    \item Request cost = 100 × \$0.0000002 = \$0.00002
    \item \textbf{Total = \$0.000145}
\end{itemize}

\pagebreak

% ==============================================================================
% SECTION 5: DATA MODELS
% ==============================================================================
\section{Data Models}

\subsection{ParsedEvent Class}
\lstset{language=Python, caption=ParsedEvent Data Class}
\begin{lstlisting}
@dataclass
class ParsedEvent:
    log_group: str
    log_stream: str
    timestamp_ms: int
    message: str
    event_type: str  # "error" | "report" | "start" | "end" | "other"
    
    # REPORT fields
    request_id: Optional[str] = None
    duration_ms: Optional[float] = None
    billed_duration_ms: Optional[float] = None
    memory_size_mb: Optional[float] = None
    max_memory_used_mb: Optional[float] = None
    init_duration_ms: Optional[float] = None
    is_cold_start: bool = False
    
    # Error fields
    error_type: Optional[str] = None
    error_message: Optional[str] = None
\end{lstlisting}

\subsection{AnalysisSummary Class}
\lstset{language=Python, caption=AnalysisSummary Data Class}
\begin{lstlisting}
@dataclass
class AnalysisSummary:
    log_groups_scanned: list[str]
    total_events_processed: int
    total_errors: int
    total_invocations: int
    cold_starts: int
    
    durations_ms: list[float]
    billed_durations_ms: list[float]
    memory_used_mb: list[float]
    memory_sizes_mb: list[float]
    error_events: list[ParsedEvent]
    estimated_cost: float
    
    # Computed properties
    avg_duration_ms, max_duration_ms, p95_duration_ms
    avg_memory_mb, memory_efficiency_pct
    cold_start_rate, error_rate
\end{lstlisting}

\pagebreak

% ==============================================================================
% SECTION 6: ERROR DETECTION PATTERNS
% ==============================================================================
\section{Error Detection Patterns}

\subsection{Regex Patterns Used}

\begin{lstlisting}[language=Python, caption=Error Detection Patterns]
_ERROR_PATTERNS = [
    re.compile(r"\[ERROR\]", re.IGNORECASE),
    re.compile(r"\bERROR\b"),
    re.compile(r"Traceback \(most recent call last\)"),
    re.compile(r"Exception:"),
    re.compile(r"Error:"),
    re.compile(r"FATAL", re.IGNORECASE),
    re.compile(r"Unhandled exception", re.IGNORECASE),
    re.compile(r"Task timed out"),
    re.compile(r"Runtime.ExitError"),
    re.compile(r"OutOfMemoryError", re.IGNORECASE),
]

_EXCEPTION_TYPE_PATTERN = re.compile(
    r"(?P<exc_type>[A-Za-z][A-Za-z0-9_]*(?:Error|Exception|Fault|Warning))"
    r"(?::\s*(?P<exc_message>.+))?"
)
\end{lstlisting}

\subsection{Error Classification}
Errors are classified by:
\begin{itemize}
    \item Direct pattern matching (ERROR, Exception, etc.)
    \item Exception type extraction from message
    \item Truncation for storage (2000 char limit)
\end{itemize}

\pagebreak

% ==============================================================================
% SECTION 7: VISUALIZATION SPECIFICATIONS
% ==============================================================================
\section{Visualization Specifications}

\subsection{Plotly Charts Configuration}

\subsubsection{Duration Histogram}
\begin{itemize}
    \item Type: px.histogram
    \item Bins: 30
    \item Color: \texttt{\#4ECDC4} (teal)
    \item Add vertical line at average
\end{itemize}

\subsubsection{Memory Histogram}
\begin{itemize}
    \item Type: px.histogram
    \item Bins: 30
    \item Color: \texttt{\#FF6B6B} (coral)
\end{itemize}

\subsubsection{Timeline Chart}
\begin{itemize}
    \item Type: go.Figure with go.Scatter
    \item Mode: lines+markers
    \item Cold starts: star markers, size 15
    \item Hover mode: x unified
\end{itemize}

\pagebreak

% ==============================================================================
% SECTION 8: DEPLOYMENT GUIDE
% ==============================================================================
\section{Deployment Guide}

\subsection{Local Deployment}

\begin{enumerate}
    \item Clone repository
    \item Create virtual environment: \texttt{python -m venv .venv}
    \item Activate: \texttt{source .venv/bin/activate}
    \item Install: \texttt{pip install -r requirements.txt}
    \item Run: \texttt{streamlit run app.py}
\end{enumerate}

\subsection{AWS Configuration}

Required IAM permissions:
\begin{lstlisting}[language=json, caption=IAM Policy]
{
    "Version": "2012-10-17",
    "Statement": [{
        "Effect": "Allow",
        "Action": [
            "logs:FilterLogEvents",
            "logs:DescribeLogStreams",
            "logs:DescribeLogGroups"
        ],
        "Resource": "*"
    }]
}
\end{lstlisting}

\pagebreak

% ==============================================================================
% APPENDIX
% ==============================================================================
\section{Appendix}

\subsection{CloudWatch Log Format Examples}

\bigskip
\noindent\textbf{REPORT Line Format:}
\begin{verbatim}
REPORT RequestId: abc-123-def  Duration: 312.45 ms  
Billed Duration: 313 ms  Memory Size: 512 MB  
Max Memory Used: 128 MB  Init Duration: 234.56 ms
\end{verbatim}

\bigskip
\noindent\textbf{START Line Format:}
\begin{verbatim}
START RequestId: abc-123-def  Version: $LATEST
\end{verbatim}

\bigskip
\noindent\textbf{Error Line Examples:}
\begin{verbatim}
[ERROR] ValueError: invalid literal for int() with base 10: 'abc'
Traceback (most recent call last):
  File "lambda_function.py", line 15, handler
    return int(value)
ValueError: invalid literal for int()
\end{verbatim}

\subsection{Dependencies}
\begin{itemize}
    \item \textbf{boto3} - AWS SDK for Python
    \item \textbf{streamlit} - Web framework
    \item \textbf{pandas} - Data manipulation
    \item \textbf{plotly} - Interactive charts
    \item \textbf{reportlab} - PDF generation
\end{itemize}

\vfill
\begin{center}
\textbf{End of Technical Documentation}
\end{center}

\end{document}
